\documentclass[../../main.tex]{subfiles}
\graphicspath{{\subfix{../../images/}}}

\begin{document}

\chapter{Compilation targets}
As one of the objectives of Fdtd Lucuma is being portable we tried to compile it on every Unix\footnote{Windows is out of scope} based system that we had access to.

\section{Linux (Debian Experimental)}
C++20's modules have not been widely implemented for a substantial amount of time as of 2026.
Thus, a bleeding edge compiler and toolchain is needed.
\inputminted[bgcolor=codeBg, tabsize=4, breaklines=true, linenos=true]{bash}{\subfix{build-debian.sh}}
\section{Linux (Arch Linux)}
Arch Linux as a rolling release distribution has no problems providing the latest software.
In addition, the Arch Linux User Repository (AUR) has niche packages.
Even if a software is not packaged, you can just package it yourself.
As of 2026 I am the maintainer of Arch Linux's Shader Slang package\footnote{\url{https://aur.archlinux.org/packages/shader-slang}}.
\inputminted[bgcolor=codeBg, tabsize=4, breaklines=true, linenos=true]{bash}{\subfix{build-arch.sh}}
\section{Linux (any other x86 distribution)}
This target primary purpose is to be run on UTEC's HPC cluster Khipu.
As student users do not have the permission to install new packages and installing all Fdtd Lucuma's dependencies
manually would be a time intensive task.
Junest\cite{junest} solves that problem by creating a nested Arch Linux environment that does not need root user privileges to use.
Unlike bare-metal Arch Linux, inside this new environment the version of the GPU driver must be the same version as the host one for the Vulkan acceleration to work.
\inputminted[bgcolor=codeBg, tabsize=4, breaklines=true, linenos=true]{bash}{\subfix{build-junest.sh}}
\section{MacOS}
MacOS does not support Vulkan natively.
Instead, it has a layer over Metal named MoltenVK\cite{moltenvk}.
\inputminted[bgcolor=codeBg, tabsize=4, breaklines=true, linenos=true]{bash}{\subfix{build-macos.sh}}
\section{Android (Unrooted Termux)}
As of 2025 most Android devices still only support up to Vulkan 1.1\cite{android_vulkan}.
Some extensions required by Ftdd Lucuma are not widely available, which forces it to use fallbacks that are not as performant.
Therefore this target tends to be slower than any other target that has an up to date version of Vulkan.
\inputminted[bgcolor=codeBg, tabsize=4, breaklines=true, linenos=true]{bash}{\subfix{build-termux.sh}}

\end{document}
